% Generated draft: requires mathematical and source review.
\documentclass[11pt]{article}
\usepackage[margin=25mm]{geometry}
\usepackage{fontspec}
\setmainfont{Latin Modern Roman}
\usepackage{amsmath,amssymb,mathtools,unicode-math}
\setmathfont{Latin Modern Math}
\usepackage{enumitem,needspace,xurl,hyperref}
\setlength{\parindent}{0pt}
\setlength{\parskip}{5pt}
\setlength{\emergencystretch}{3em}
\newcommand{\N}{\mathbb{N}}
\newcommand{\Z}{\mathbb{Z}}
\newcommand{\Q}{\mathbb{Q}}
\newcommand{\R}{\mathbb{R}}
\newcommand{\C}{\mathbb{C}}
\newcommand{\F}{\mathbb{F}}
\newcommand{\A}{\mathbb{A}}
\newcommand{\PP}{\mathbb{P}}
\newcommand{\E}{\mathbb{E}}
\newcommand{\eps}{\varepsilon}
\newcommand{\dd}{\,\mathrm{d}}
\begin{document}
\section*{UnsolvedMath ID: 30000970}
\textbf{Source problem code:} catalogue entry
\subsection{Definitions and mathematical statement}
Let \(\mathcal X\) be an Artin stack over an unspecified base. In the usual definition, a good moduli space is a morphism \(\phi:\mathcal X\to X\) to an algebraic space such that \(\phi\) is cohomologically affine and the natural map \(\mathcal O_X\to\phi_*\mathcal O_{\mathcal X}\) is an isomorphism. Here cohomologically affine means quasi-compact with exact pushforward on quasi-coherent modules.

The source poses three directions for good moduli spaces:
\begin{enumerate}
\item Is there a Hilbert--Mumford-type criterion in this setting? The numerical data, degeneration tests, and precise property to be characterized are not specified.
\item Do suitable topological criteria on an Artin stack, such as an intended weak valuative criterion, imply existence of a good moduli space? The task is to formulate effective sufficient conditions; the extraction supplies no particular weak valuative condition or additional stack hypotheses.
\item Can good moduli spaces be constructed locally? The intended topology or system of local charts, compatibility of local constructions, and gluing statement are not specified.
\end{enumerate}
These are exploratory existence and construction questions. They do not assert that every Artin stack has a good moduli space or that an unnamed valuative condition alone is sufficient.
\subsection{Short English statement}
Seek a Hilbert--Mumford-type criterion, topological or weak valuative existence criteria, and local construction methods for good moduli spaces of Artin stacks. The source does not fix the hypotheses or the exact criteria to be tested.
\subsection{Sources}
\begin{itemize}
\item Local UnsolvedMath registry, permanent ID \(30000970\), source code \texttt{catalogue entry}, ``Existence Criteria for Good Moduli Spaces.'' Registry curation: Ulam AI / UnsolvedMath Contributors, CC BY 4.0; underlying sources retain their supplied provenance.
\item \emph{Classical Algebraic Geometry} (2009), Oberwolfach Reports, pp.\ 1425--1474, DOI \url{https://doi.org/10.4171/owr/2008/26}.
\item Supplied ``Alper, local properties of good moduli spaces'': \url{https://doi.org/10.2748/tmj/1332767342}.
\item Supplied ``Alper--Fedorchuk--Smyth existence criterion'': \url{https://doi.org/10.1112/S0010437X16008289}.
\item The supplied assessment checked 2026-08-21 labels the item partially solved, saying local structure, criteria and constructions are known in important settings but that no unrestricted solution to all three questions exists. This is dated curation, not independent verification or a claim that every special case remains open.
\end{itemize}
\end{document}
